% !TEX root = ../main.tex
\clearpage
\phantomsection
\section*{3 \textbf{Literature Review}}
\label{sec:literature-review}
\addcontentsline{toc}{section}{3 Literature Review}

\phantomsection
\subsection*{3.1 Graph-Based On-Chain Reputation Scoring}
\label{sec:graph-based-on-chain-reputation-scoring}
\addcontentsline{toc}{subsection}{3.1 Graph-Based On-Chain Reputation Scoring}

Although prior work often uses Ethereum mainnet as the reference chain, \textbf{this dissertation's primary empirical setting is Arbitrum One with GMX V2 perpetual swap position outcomes} for domain validation. The literature below remains relevant as background for graph-based reputation methods.

Early research on on-chain reputation leveraged graph theory, particularly adaptations of PageRank, to rank wallet addresses by ``trustworthiness'' (Page et al., 1998). Do et al. (2023) model Ethereum addresses as nodes in a directed graph with transactions as edges, then apply the PageRank algorithm to compute a social credit score for each address. The rationale is that DeFi's growth demands methods to identify and rank entities for unsecured lending and customer protection. In this framework, an address accumulates higher score if it is ``trusted'' by many or large transactions from others, analogous to how important web pages have many incoming links. PageRank thus measures connectivity strength and activity as proxies for reputation. Subsequent graph models combine multiple factors (e.g., frequency and volume of interactions) into edge weights (Do et al., 2023; Vathsala V. \& V. D. Sharma, 2025).


A key consideration in graph construction is how to weight the edges. Basic implementations treat each transaction equally, but later work incorporates transaction value as a weight, so that large transfers signify stronger trust. Using a HodgeRank approach, Do et al. (2023) treat the amount of ETH transferred as the weight on each edge and then solve for a global ranking that accounts for these weighted flows. This produces an Ethereum reputation ranking influenced by transaction volume, under the intuition that high-value interactions indicate greater trust or stake between parties. Another practical challenge is dealing with graph structure: Ethereum's transaction graph can contain many isolated clusters or "sink" nodes that absorb PageRank score. Damping factors, as in traditional PageRank, are used to mitigate score inflation for sinks and small isolated subgraphs, ensuring the reputation scores reflect participation in the broader network rather than siloed activity.


One limitation of a vanilla PageRank, i.e., the unmodified PageRank formulation applied to a transaction graph without explicit risk propagation or domain-specific adjustments, is that it may not distinguish malicious but active actors from genuinely reputable ones. Lin et al. (2024) show that transaction connectivity alone can miss risk-relevant structure unless risk signals are propagated through the network. This highlights the need for refined models. Proposed improvements include incorporating negative feedback or propagation of risk labels through the network. For instance, recent research introduced RiskProp, a network-propagation-based risk scoring model for Ethereum accounts (Lin et al., 2024). RiskProp augments the transaction graph with a "de-anonymous score" for addresses and propagates risk signals through directed bipartite graphs. The approach was able to flag a majority of known high-risk accounts by propagating risk indicators across transaction links, effectively identifying suspicious wallets that a simple connectivity measure might overlook. This illustrates how graph-based methods can be extended beyond pure PageRank to serve as on-chain analogues of credit risk models, where the "credit" in question may be financial reliability or compliance risk.


It is worth noting that graph-based reputation systems often consider temporal dynamics. Real-world creditworthiness changes over time, so some models incorporate time-decay factors to diminish the influence of stale interactions. Do et al.\ (2023) apply logistic time-decay to transfers in AWP so that recent activity contributes more strongly than stale edges, while Lin et al.\ (2024) show that risk-relevant structure can be masked unless propagation models are updated as network state evolves. A dynamic reputation model might weight recent transactions more heavily or periodically expire old edges. This prevents entities from resting on long-past positive interactions and addresses the whitewashing problem, where an entity's bad behavior gets diluted by time. However, striking the right balance is crucial: too much decay can ignore longstanding trust relationships, while too little can allow outdated information to skew the scores.


\phantomsection
\subsection*{3.2 Predictive Modeling for On-Chain Credit Risk}
\label{sec:predictive-modeling-for-on-chain-credit-risk}
\addcontentsline{toc}{subsection}{3.2 Predictive Modeling for On-Chain Credit Risk}

In parallel with graph-theoretical methods, researchers have developed data-driven machine-learning models for on-chain credit scoring, especially in DeFi lending contexts. These models take inspiration from traditional credit analytics, using wallet histories as input features and predicting the likelihood of default or delinquency. Ghosh et al.\ (2024), in an arXiv preprint, propose an on-chain credit risk scoring framework for decentralized finance that featurizes wallet activity from public event logs and trains supervised models to rank borrowers by repayment and liquidation risk. In their framework, each wallet's activity is featurized analogously to a credit report: features include payment history, amounts owed, length of credit history, new credit, and credit mix, all derived from on-chain behavior. Using machine-learning classifiers, the model estimates the probability that a given on-chain lending position will become delinquent (e.g., fall below collateral requirements and face liquidation) in the near future. The output is essentially a credit risk score for the wallet or position, indicating the odds of default. Ghosh et al.\ report that this approach can successfully rank borrowers by risk and thus support under-collateralized lending by relying on predicted probability of repayment rather than collateral ratios alone.


Other machine-learning approaches similarly treat on-chain transactions and account data as inputs to predictive models for creditworthiness. Hassija et al. (2020) presented a blockchain-based credit scoring model that leverages prospect theory to weight positive vs. negative credit events differently, capturing the asymmetric risk attitudes of borrowers. Their model, implemented as a smart contract, calculates a trust score for each participant which is updated as new lending outcomes occur.


By incorporating prospect theory (which, for instance, might penalize defaults more heavily than it rewards timely repayments), the scoring mechanism aligns with human risk biases, potentially improving prediction of who will default. Similarly, Chakraborty et al.\ (2019) propose a blockchain-based credit analysis framework that stores and evaluates customer credit data on-chain for more efficient financial-system workflows. Their design demonstrates how distributed ledger structures can support auditable credit assessment alongside conventional transaction histories.


\phantomsection
\subsection*{3.3 Industry Initiatives and Scholarly Perspectives on DeFi Credit Scoring}
\label{sec:industry-initiatives-and-non-scholarly-perspectives}
\addcontentsline{toc}{subsection}{3.3 Industry Initiatives and Scholarly Perspectives on DeFi Credit Scoring}

Industry demand for on-chain credit scoring has prompted both commercial deployments and scholarly analysis of decentralized lending workflows. Moln\'{a}r et al.\ (2023) examine blockchain-based business process management in finance and show how credit and claim requests can be re-engineered through smart-contract workflows, highlighting the need for transparent, auditable risk signals in DeFi lending. Packin and Lev-Aretz (2024) analyze hybrid on-chain/off-chain credit architectures and the governance challenges of integrating traditional credit data with blockchain-native activity, underscoring that reputation metrics must remain interpretable under pseudonymity constraints.


Recent scholarly work also treats DeFi credit scoring as a measurable research problem rather than a purely commercial feature. Shimpi (2024) implements an Ethereum-based lending workflow that integrates traditional credit-scoring components with smart-contract credit bureau functions, illustrating how off-chain credit data can be bridged to on-chain lending decisions. Vathsala V. and V. D. Sharma (2025) propose a secure AI architecture for dynamic DeFi credit scoring that combines federated learning with blockchain-anchored data provenance, noting that most deployed systems rely on transfer histories or loan outcomes rather than explicit authorization relationships.


These scholarly perspectives highlight current gaps that EndorseRank aims to fill. Existing implementations focus heavily on past transaction behavior (payments, liquidations, balances) or supervised default prediction, but they do not systematically leverage ERC-20 token approval relationships between wallets as primary graph edges. An ERC-20 approve (or permit) action, which authorizes another address to spend tokens on the owner's behalf (Ethereum Improvement Proposals, 2015; Ethereum Improvement Proposals, 2020), is a strong signal of trust or financial linkage that lies outside direct token transfers. Current graph-based scoring systems largely ignore these approval edges, focusing only on realized transfers of value. EndorseRank addresses this research gap by integrating ERC-20 approve/permit edges as first-class links in the reputation graph, treating them as directed weighted edges in a PageRank computation. Because trust in operational DeFi systems is established not only through the scholarly models reviewed so far but also through widely deployed identity and reputation tooling, Section~3.4 examines that practice explicitly and states how EndorseRank differs from it.


Furthermore, EndorseRank is distinctive in its handling of temporal effects. Whereas many reputation systems apply time-based decay to down-weight old interactions, EndorseRank proposes to forgo time decay for approval/permit edges. The rationale is that an approval remains in force until revoked; a permission granted a year ago is equally pertinent today if the allowance is still active. By giving persistent weight to long-lived trust links, EndorseRank preserves important information about enduring financial relationships that a short memory window would lose. This contrasts with typical transaction-based scoring, where a payment from years ago might reasonably be given little importance. In the context of approvals, the lack of time decay reflects the enduring nature of the risk: as long as B is authorized to spend A's funds, A's exposure (and hence the reputational linkage) does not diminish over time. This no-decay approach is distinctive, as most prior works have either not included such edges at all or have not considered the temporal aspect of trust explicitly. By integrating ERC-20 approval networks with PageRank and maintaining their influence over time, EndorseRank is expected to provide a richer and more stable on-chain reputation score---one that captures both the flow of value and the web of delegated permissions on Arbitrum One.


\phantomsection
\subsection*{3.4 Attestation-Based Reputation Tooling in Deployed DeFi Systems}
\label{sec:attestation-based-reputation-tooling}
\addcontentsline{toc}{subsection}{3.4 Attestation-Based Reputation Tooling in Deployed DeFi Systems}

Alongside the scholarly models reviewed above, wallet-level trust in operational DeFi and web3 systems is most commonly established through identity and reputation tooling bound directly to a wallet address. The most widely deployed instance is Human Passport (formerly Gitcoin Passport), which aggregates verifiable credentials---termed Stamps---issued for government-identity (KYC) checks, biometrics, web-of-trust vouching, and prior web2 and web3 activity. Each Stamp carries a weight that reflects its cost of forgery, a wallet's score is the weighted sum of the Stamps it holds, and that score is compared against a passing threshold used to admit or exclude the address (Human Passport, n.d.). Stamps are verified off-chain and may optionally be minted on-chain as attestations through the Ethereum Attestation Service, so that smart contracts can read a wallet's score without relying on an external server (Human Passport, n.d.).


Two properties of this design bear directly on the present study. First, the score deliberately combines on-chain and off-chain evidence, because its purpose is Sybil resistance: establishing that an address is controlled by a distinct person rather than being one of many addresses operated by a single actor. Second, the score functions largely as a gate---a wallet either meets the threshold or it does not---rather than as a graded ordering of wallets by financial reliability. Comparative security analysis classifies Passport within the attestation-stacking family of Sybil defences, reports strong but not decisive discrimination relative to biometric and social-graph alternatives, and identifies credential farming, in which an adversary acquires low-cost Stamps in bulk, as its principal residual weakness (Bordeianu \& Popescu, 2026). The tooling originates in quadratic-funding mechanisms, whose matching allocations rest on a one-contributor-one-voice assumption and are therefore acutely sensitive to duplicated identities (Buterin et al., 2019).


This practice reflects a broader line of work on proof of personhood. Siddarth et al.\ (2020) review protocols that establish Sybil-resistant identity through subjective inputs such as voting, vouching, and interpretation, and observe that such systems must trade off Sybil resistance, self-sovereignty, and privacy against one another. Whatever substrate is used, the question these protocols answer is whether an identifier corresponds to a unique person, not how one wallet should be ranked against another for credit-like purposes.


Passport additionally offers model-based scoring that classifies any EVM address as human or Sybil in real time from its on-chain activity across Ethereum and several Layer-2 networks, including Arbitrum, by comparing that activity against curated lists of known Sybil and human addresses (Human Passport, n.d.). This confirms that purely on-chain behavior carries usable trust signal on the same chain studied here. The features involved, however, are per-address activity aggregates evaluated against a binary human-or-Sybil target; the directed relationships between wallets are not themselves the object of measurement.


EndorseRank is therefore neither a competitor to nor a replacement for attestation-based tooling, and the difference is best stated along five dimensions:


\begin{itemize}
\item[\textbullet] \textbf{Signal source}: Passport aggregates third-party attestations about the person behind an address; EndorseRank derives its edges from a financial act that wallets themselves record on-chain, namely an ERC-20 approval authorizing a spender to move the owner's tokens.
\item[\textbullet] \textbf{Data scope}: Passport combines off-chain credentials (KYC, biometrics, web2 accounts) with on-chain activity; EndorseRank uses public on-chain state only and requires no personal data, preserving permissionless participation.
\item[\textbullet] \textbf{Question answered}: Passport asks whether an address is a distinct human; EndorseRank asks how a wallet ranks relative to other wallets within a network of delegated spending authority.
\item[\textbullet] \textbf{Output form}: Passport produces a score interpreted against a pass threshold; EndorseRank produces a continuous, network-propagated ranking suitable for differentiated risk parameters.
\item[\textbullet] \textbf{Validation target}: Passport-style scores are assessed against Sybil-labelled address sets in grant and airdrop distribution; EndorseRank is assessed against realized GMX V2 perpetual trading outcomes on Arbitrum One.
\end{itemize}
The two are complementary rather than mutually exclusive: a personhood gate can remove duplicated identities before a relational reputation score is computed over the remaining population, and Section~6 treats the absence of such a gate in this study as an explicit limitation. What neither the scholarly literature of Sections~3.1--3.3 nor the deployed tooling described here has examined is whether the delegation relationships already recorded in ERC-20 allowances constitute a usable trust signal in their own right. That is the specific gap this study addresses, and it is the basis on which the study claims an original contribution.


In summary, the literature shows a progression from simple graph-based reputation metrics to machine-learning models for on-chain credit scoring, running in parallel with an industry practice that scores wallets through aggregated identity attestations. PageRank-based techniques contribute a transparent, explainable, network-centric view of trust and influence; machine-learning models offer risk-ranking performance using diverse features; attestation-based tooling supplies personhood assurance but not a graded ordering of financial reliability. A gap remains between them: incorporating meaningful relational graph features, such as ERC-20 approval links, into a scoring system that is at once explainable, permissionless, and domain-aligned. EndorseRank aims to bridge that gap. It builds on the insight that network structure matters---who authorizes whom---and that this structure can be quantified through a PageRank-style algorithm. Empirical validation in this dissertation uses GMX V2 perpetual swap position outcomes on Arbitrum One: close-success count and realized gain proxy serve as observable trading-success measures, analogous to the in-degree and in-value proxies used in Do et al.\ (2023). By doing so without arbitrarily discarding older trust links, EndorseRank preserves the long-term context of relationships. This approach addresses limitations in prior art and aligns with scholarly calls for transparent, composable DeFi risk assessment (Moln\'{a}r et al., 2023; Shimpi, 2024). It promises an enhanced on-chain reputation score that could improve risk assessment for leveraged DeFi activity on Arbitrum One, making decentralized finance both more accessible and more secure.
